% ==============================================================================
%  รายงานโครงงาน: Magic Spell Duel
%  วิธีคอมไพล์: ใช้ XeLaTeX (บน Overleaf: Menu > Compiler > XeLaTeX) และคอมไพล์ 2 รอบเพื่อให้สารบัญสมบูรณ์
%  ฟอนต์: เลือกอัตโนมัติจากฟอนต์ภาษาไทยที่มีในเครื่อง (TH Sarabun New > Angsana New > ฟอนต์ชุด TLWG > Noto)
% ==============================================================================
\documentclass[12pt,a4paper]{report}

\usepackage[top=2.5cm,bottom=2.5cm,left=3.2cm,right=2.5cm]{geometry}
\usepackage{fontspec}
\usepackage{amsmath}
\usepackage{graphicx}
\usepackage{booktabs}
\usepackage{tabularx}
\usepackage{array}
\usepackage{enumitem}
\usepackage{xcolor}
\usepackage{listings}
\usepackage{float}
\usepackage{indentfirst}
\usepackage{tikz}
\usetikzlibrary{arrows.meta,positioning}
\usepackage[hidelinks,unicode]{hyperref}

% ---------- ฟอนต์ภาษาไทย ----------
\XeTeXlinebreaklocale "th"
\XeTeXlinebreakskip = 0pt plus 1pt
\IfFontExistsTF{TH Sarabun New}{%
  \setmainfont[Scale=1.33]{TH Sarabun New}%
}{\IfFontExistsTF{Angsana New}{%
  \setmainfont[Scale=1.33]{Angsana New}%
}{\IfFontExistsTF{Laksaman}{%
  \setmainfont{Laksaman}%
}{\IfFontExistsTF{Garuda}{%
  \setmainfont{Garuda}%
}{\IfFontExistsTF{Noto Serif Thai}{%
  \setmainfont{Noto Serif Thai}%
}{%
  \setmainfont{Noto Sans Thai}%
}}}}}
\linespread{1.15}
\setlength{\parindent}{1.2cm}
\setlength{\emergencystretch}{3em}

% ---------- ชื่อหัวข้อภาษาไทย ----------
\renewcommand{\contentsname}{สารบัญ}
\renewcommand{\listfigurename}{สารบัญรูป}
\renewcommand{\listtablename}{สารบัญตาราง}
\renewcommand{\chaptername}{บทที่}
\renewcommand{\appendixname}{ภาคผนวก}
\renewcommand{\figurename}{รูปที่}
\renewcommand{\tablename}{ตารางที่}
\renewcommand{\bibname}{บรรณานุกรม}
\renewcommand{\lstlistingname}{โค้ดที่}

% ---------- คำสั่งช่วย ----------
\newcommand{\code}[1]{\texttt{\detokenize{#1}}}
\newcolumntype{Y}{>{\raggedright\arraybackslash}X}

\lstset{
  language=Python,
  basicstyle=\ttfamily\footnotesize,
  keywordstyle=\color{blue!60!black},
  commentstyle=\color{green!40!black},
  stringstyle=\color{red!55!black},
  numbers=left,
  numberstyle=\tiny\color{gray},
  frame=single,
  rulecolor=\color{gray!60},
  breaklines=true,
  showstringspaces=false,
  columns=fullflexible,
  keepspaces=true,
  xleftmargin=1.5em
}

\tikzset{
  box/.style={draw, rounded corners=2pt, align=center, minimum height=0.9cm, inner sep=5pt, fill=gray!8},
  thr/.style={draw, rounded corners=2pt, align=center, minimum height=0.9cm, inner sep=5pt, fill=blue!8},
  state/.style={draw, rounded corners=8pt, align=center, minimum height=0.9cm, minimum width=2.4cm, fill=orange!12},
  arr/.style={-{Stealth[length=2.2mm]}, thick}
}

% ==============================================================================
%  ข้อมูลหน้าปก  <<< แก้ไขข้อมูลในส่วนนี้ก่อนส่ง >>>
% ==============================================================================
\newcommand{\ReportAuthor}{[ชื่อ--นามสกุล ผู้จัดทำ]}
\newcommand{\ReportStudentID}{[รหัสนักศึกษา]}
\newcommand{\ReportCourse}{[รหัสวิชา และชื่อรายวิชา]}
\newcommand{\ReportAdvisor}{[ชื่ออาจารย์ผู้สอน]}
\newcommand{\ReportFaculty}{[สาขาวิชา / คณะ]}
\newcommand{\ReportUniversity}{[ชื่อมหาวิทยาลัย]}
\newcommand{\ReportTerm}{[ภาคเรียนที่ ... ปีการศึกษา ...]}

\begin{document}

% ==============================================================================
%  หน้าปก
% ==============================================================================
\begin{titlepage}
  \centering
  \vspace*{0.5cm}
  \IfFileExists{../assets/images/title_logo.png}{%
    \includegraphics[width=0.62\textwidth]{../assets/images/title_logo.png}\par
  }{}
  \vspace{1.2cm}
  {\Huge\bfseries รายงานโครงงาน\par}
  \vspace{0.6cm}
  {\LARGE\bfseries Magic Spell Duel\par}
  \vspace{0.3cm}
  {\Large เกมต่อสู้เวทมนตร์สำหรับผู้เล่นสองคน\\ควบคุมด้วยท่าทางมือผ่านกล้องเว็บแคม\par}
  \vfill
  {\Large จัดทำโดย\par}
  \vspace{0.2cm}
  {\Large \ReportAuthor\par}
  {\Large รหัสนักศึกษา \ReportStudentID\par}
  \vfill
  {\Large เสนอ\par}
  \vspace{0.2cm}
  {\Large \ReportAdvisor\par}
  \vfill
  {\large รายงานนี้เป็นส่วนหนึ่งของรายวิชา \ReportCourse\par}
  {\large \ReportFaculty\par}
  {\large \ReportUniversity\par}
  {\large \ReportTerm\par}
\end{titlepage}

\pagenumbering{roman}

% ==============================================================================
%  บทคัดย่อ
% ==============================================================================
\chapter*{บทคัดย่อ}
\addcontentsline{toc}{chapter}{บทคัดย่อ}

โครงงานนี้พัฒนาเกม \textbf{Magic Spell Duel} ซึ่งเป็นเกมต่อสู้สำหรับผู้เล่นสองคนที่ยืนอยู่หน้ากล้องเว็บแคมตัวเดียวกัน
ผู้เล่นไม่ต้องใช้คีย์บอร์ดหรือจอยควบคุม แต่ร่ายคาถาด้วยการ ``วาดสัญลักษณ์กลางอากาศ'' ด้วยนิ้วชี้
แล้วทำท่ามือที่กำหนดเพื่อปล่อยคาถาใส่คู่ต่อสู้ ระบบพัฒนาด้วยภาษา Python โดยใช้ไลบรารี OpenCV
สำหรับรับภาพและวาดกราฟิก MediaPipe สำหรับตรวจจับจุดสำคัญของมือ 21 จุดและตรวจจับใบหน้า
และ Pygame สำหรับระบบเสียง

แกนหลักของระบบประกอบด้วย (1) การจำแนกท่ามือด้วยกฎเชิงเรขาคณิตบนจุดสำคัญของมือ
(2) การรู้จำลายเส้นด้วยอัลกอริทึม \$1 Unistroke Recognizer ที่ดัดแปลงให้รองรับการวาดย้อนทิศและการเริ่มวาดจากจุดใดก็ได้
(3) เครื่องสถานะ (State Machine) ที่ควบคุมวงจรการร่ายคาถา มานา คูลดาวน์ และกติกาการแข่งขันแบบชนะ 2 ใน 3 ยก
และ (4) ระบบเอฟเฟกต์ภาพ เสียง และส่วนแสดงสถานะบนจอ

นอกจากนี้ได้วิเคราะห์และปรับปรุงประสิทธิภาพของโปรแกรม โดยแยกการอ่านภาพจากกล้องและการประมวลผลโมเดล
ออกไปทำงานบนเธรดของตนเอง ปรับวิธีเบลนด์ภาพเอฟเฟกต์ และแปลงตัวรู้จำลายเส้นให้คำนวณแบบเวกเตอร์
ผลการวัดบนเครื่องที่ใช้พัฒนาพบว่าอัตราการแสดงผลเพิ่มจากประมาณ 17--30 เฟรมต่อวินาที เป็นประมาณ 85--88 เฟรมต่อวินาที
ในสภาวะปกติ อัตราเฟรมของกล้องเพิ่มจาก 30 เป็นประมาณ 60 เฟรมต่อวินาที และระยะเวลาที่เกมหยุดชะงักเมื่อรู้จำลายเส้น
ลดจากประมาณ 270 มิลลิวินาทีเหลือประมาณ 8 มิลลิวินาที อีกทั้งยังปรับให้ตรวจจับมือที่มีขนาดเล็กลงประมาณครึ่งหนึ่งของเดิมได้
เพื่อรองรับผู้เล่นที่ยืนห่างจากกล้อง

\vspace{0.4cm}
\noindent\textbf{คำสำคัญ:} การรู้จำท่าทางมือ, คอมพิวเตอร์วิทัศน์, MediaPipe, \$1 Unistroke Recognizer, เกมควบคุมด้วยการเคลื่อนไหว

\tableofcontents
\listoffigures
\listoftables

\clearpage
\pagenumbering{arabic}

% ==============================================================================
\chapter{บทนำ}
% ==============================================================================

\section{ที่มาและความสำคัญ}
เกมที่ควบคุมด้วยการเคลื่อนไหวของร่างกายเคยต้องอาศัยอุปกรณ์เฉพาะทาง เช่น กล้องวัดความลึกหรือจอยที่มีเซนเซอร์
แต่ปัจจุบันโมเดลการเรียนรู้ของเครื่องสำหรับตรวจจับมือสามารถทำงานแบบเรียลไทม์ได้บนหน่วยประมวลผลกลางของคอมพิวเตอร์ทั่วไป
โดยใช้เพียงกล้องเว็บแคมธรรมดา จึงเปิดโอกาสให้สร้างเกมลักษณะนี้ได้โดยไม่ต้องมีฮาร์ดแวร์เพิ่มเติม

โครงงานนี้จึงนำเทคนิคด้านคอมพิวเตอร์วิทัศน์มาประยุกต์เป็นเกมต่อสู้เวทมนตร์ ซึ่งผู้เล่นต้อง ``วาด'' สัญลักษณ์ของคาถาในอากาศ
และ ``ทำท่ามือ'' เพื่อปล่อยคาถา การออกแบบเช่นนี้ทำให้ผู้จัดทำได้ฝึกทั้งการประมวลผลภาพ การออกแบบอัลกอริทึมรู้จำรูปแบบ
การจัดการสถานะของโปรแกรมแบบเรียลไทม์ และการวิเคราะห์ประสิทธิภาพของซอฟต์แวร์

\section{วัตถุประสงค์}
\begin{enumerate}[leftmargin=2cm]
  \item เพื่อพัฒนาเกมสำหรับผู้เล่นสองคนที่ควบคุมด้วยท่าทางมือผ่านกล้องเว็บแคมเพียงตัวเดียว
  \item เพื่อศึกษาและประยุกต์ใช้การตรวจจับจุดสำคัญของมือ (Hand Landmark Detection) ในการจำแนกท่ามือ
  \item เพื่อศึกษาและพัฒนาการรู้จำลายเส้นที่วาดกลางอากาศด้วยอัลกอริทึม \$1 Unistroke Recognizer
  \item เพื่อออกแบบระบบเกมที่มีกติกา สมดุล เอฟเฟกต์ภาพและเสียงครบถ้วน และทำงานได้ลื่นไหลแบบเรียลไทม์
\end{enumerate}

\section{ขอบเขตของโครงงาน}
\begin{itemize}[leftmargin=2cm]
  \item รองรับผู้เล่น 2 คน โดยแบ่งแดนซ้าย (ผู้เล่น 1) และขวา (ผู้เล่น 2) ของภาพจากกล้อง
  \item ใช้กล้องเว็บแคม 1 ตัว ที่ความละเอียด 1280$\times$720 พิกเซล บนระบบปฏิบัติการ Windows
  \item มีคาถาพื้นฐาน 5 ชนิด และคาถาพิเศษ 1 ชนิดที่ต้องปลดล็อกระหว่างการแข่งขัน
  \item แข่งขันแบบชนะ 2 ใน 3 ยก (Best of 3)
  \item ประมวลผลบนหน่วยประมวลผลกลาง (CPU) ทั้งหมด ไม่ใช้หน่วยประมวลผลกราฟิก
\end{itemize}

\section{ประโยชน์ที่คาดว่าจะได้รับ}
\begin{itemize}[leftmargin=2cm]
  \item ได้เกมต้นแบบที่เล่นได้จริงโดยไม่ต้องใช้อุปกรณ์ควบคุมเพิ่มเติม
  \item ได้ความรู้และประสบการณ์ในการนำโมเดลสำเร็จรูปมาประกอบเป็นระบบแบบเรียลไทม์
  \item ได้แนวทางการวัดและปรับปรุงประสิทธิภาพของโปรแกรมที่มีงานประมวลผลหนัก
\end{itemize}

\section{เครื่องมือที่ใช้ในการพัฒนา}
\begin{table}[H]
  \centering
  \caption{เครื่องมือและไลบรารีที่ใช้}
  \label{tab:tools}
  \begin{tabularx}{\textwidth}{l Y}
    \toprule
    \textbf{เครื่องมือ} & \textbf{หน้าที่ในโครงงาน} \\
    \midrule
    Python 3.12 & ภาษาหลักที่ใช้พัฒนา \\
    OpenCV & รับภาพจากกล้อง แปลงสี วาดกราฟิกและข้อความ แสดงผลหน้าต่างเกม \\
    MediaPipe & โมเดลตรวจจับจุดสำคัญของมือ 21 จุด และโมเดลตรวจจับใบหน้า \\
    NumPy & คำนวณเชิงเวกเตอร์สำหรับการเบลนด์ภาพและการรู้จำลายเส้น \\
    Pygame (mixer) & เล่นเพลงประกอบ เสียงเอฟเฟกต์ และเสียงผู้ประกาศ \\
    Pillow & สร้างภาพไอคอนอีโมจิสำหรับคำแนะนำท่ามือบนจอ \\
    Git / GitHub & ควบคุมเวอร์ชันของซอร์สโค้ด \\
    \bottomrule
  \end{tabularx}
\end{table}

% ==============================================================================
\chapter{ทฤษฎีและเทคโนโลยีที่เกี่ยวข้อง}
% ==============================================================================

\section{การตรวจจับจุดสำคัญของมือด้วย MediaPipe Hands}
MediaPipe Hands \cite{zhang2020hands} ทำงานเป็นสองขั้นตอน ขั้นแรกคือ \textit{โมเดลตรวจจับฝ่ามือ} (Palm Detector)
ซึ่งค้นหาตำแหน่งของฝ่ามือจากภาพทั้งภาพ ขั้นที่สองคือ \textit{โมเดลจุดสำคัญของมือ} (Hand Landmark Model)
ซึ่งรับภาพที่ตัดเฉพาะบริเวณมือแล้วทำนายพิกัดของจุดสำคัญ 21 จุด ได้แก่ ข้อมือ 1 จุด และข้อนิ้วของแต่ละนิ้วนิ้วละ 4 จุด
เมื่อตรวจพบมือแล้ว เฟรมถัดไปจะใช้ตำแหน่งมือจากเฟรมก่อนหน้าเป็นจุดเริ่มต้น จึงไม่ต้องรันโมเดลตรวจจับฝ่ามือทุกเฟรม

ข้อสังเกตที่สำคัญต่อโครงงานนี้คือ โมเดลตรวจจับฝ่ามือย่อภาพที่รับเข้าไปเป็นขนาดเล็กคงที่ก่อนประมวลผล
มือที่อยู่ไกลกล้องจึงเหลือรายละเอียดน้อยมากและอาจตรวจไม่พบ ประเด็นนี้นำไปสู่การปรับปรุงในหัวข้อ~\ref{sec:smallhand}

\section{การตรวจจับใบหน้าด้วย MediaPipe Face Detection}
MediaPipe Face Detection พัฒนาบนโมเดล BlazeFace \cite{bazarevsky2019blazeface} ซึ่งออกแบบให้ทำงานได้เร็วบนอุปกรณ์ทั่วไป
ผลลัพธ์คือกรอบสี่เหลี่ยมล้อมใบหน้าพร้อมค่าความเชื่อมั่น โครงงานนี้ใช้ตำแหน่งใบหน้าเพื่อประมาณตำแหน่งลำตัวของผู้เล่น
ซึ่งใช้เป็นเป้ารับความเสียหาย (Hitbox) และใช้ประมาณระยะห่างของผู้เล่นจากกล้อง

\section{อัลกอริทึม \$1 Unistroke Recognizer}
\$1 Recognizer \cite{wobbrock2007dollar} เป็นอัลกอริทึมรู้จำลายเส้นแบบเส้นเดียวที่ไม่ต้องฝึกโมเดล
ใช้เพียงการเปรียบเทียบลายเส้นของผู้ใช้กับแม่แบบ (Template) ที่เตรียมไว้ มีขั้นตอนหลักดังนี้
\begin{enumerate}[leftmargin=2cm]
  \item \textbf{Resample} เกลี่ยจุดบนลายเส้นให้มีจำนวนคงที่ $N$ จุด และห่างเท่ากันตลอดความยาวเส้น
  \item \textbf{Scale} ปรับขนาดกรอบล้อมลายเส้นให้เป็นสี่เหลี่ยมจัตุรัสขนาดมาตรฐาน
  \item \textbf{Translate} เลื่อนจุดศูนย์กลางมวลของลายเส้นมาอยู่ที่จุดกำเนิด
  \item \textbf{Match} คำนวณระยะห่างเฉลี่ยแบบจุดต่อจุดกับแม่แบบแต่ละอัน โดยค้นหามุมหมุนที่ให้ระยะห่างน้อยที่สุด
        ด้วยวิธี Golden Section Search
\end{enumerate}
ระยะห่างเฉลี่ยระหว่างลายเส้น $C$ กับแม่แบบ $T$ และคะแนนความคล้ายคำนวณได้จาก
\begin{equation}
  d(C,T) = \frac{1}{N}\sum_{i=1}^{N}\sqrt{(C_{i,x}-T_{i,x})^2 + (C_{i,y}-T_{i,y})^2},
  \qquad
  \text{score} = 1 - \frac{d_{\min}}{\tfrac{1}{2}\sqrt{S^2+S^2}}
  \label{eq:score}
\end{equation}
โดย $S$ คือความยาวด้านของกรอบมาตรฐาน และ $d_{\min}$ คือระยะห่างที่น้อยที่สุดจากแม่แบบทั้งหมด

\section{การกรองสัญญาณด้วยค่าเฉลี่ยเคลื่อนที่แบบเอกซ์โพเนนเชียล}
พิกัดที่ได้จากโมเดลมีการสั่นเล็กน้อยในแต่ละเฟรม จึงใช้ค่าเฉลี่ยเคลื่อนที่แบบเอกซ์โพเนนเชียล (Exponential Moving Average)
เพื่อลดการสั่นโดยยังตอบสนองต่อการเคลื่อนที่จริงได้
\begin{equation}
  s_t = \alpha\,x_t + (1-\alpha)\,s_{t-1}
  \label{eq:ema}
\end{equation}
โดย $x_t$ คือค่าที่วัดได้ $s_t$ คือค่าที่กรองแล้ว และ $\alpha \in (0,1]$ กำหนดน้ำหนักของค่าล่าสุด

\section{การเบลนด์ภาพ}
เอฟเฟกต์ในเกมใช้การเบลนด์สองแบบ ให้ $F$ เป็นสีของภาพเอฟเฟกต์ $B$ เป็นสีของพื้นหลัง และ $a \in [0,1]$ เป็นค่าความทึบ
\begin{align}
  \text{Alpha Blending:}\quad & O = a\,F + (1-a)\,B \\
  \text{Additive Blending:}\quad & O = \min\bigl(255,\; B + g\,a\,F\bigr)
\end{align}
แบบแรกใช้กับวัตถุทึบ เช่น ก้อนหิน ส่วนแบบที่สองใช้กับแสงเรือง เช่น ลูกไฟและสายฟ้า โดย $g$ เป็นค่าขยายความสว่าง

\section{เครื่องสถานะจำกัด}
เครื่องสถานะจำกัด (Finite State Machine) เป็นแบบจำลองที่ระบบอยู่ในสถานะใดสถานะหนึ่งเสมอ
และเปลี่ยนสถานะเมื่อเกิดเหตุการณ์ที่กำหนด โครงงานนี้ใช้เครื่องสถานะสองระดับ คือระดับการแข่งขันและระดับผู้เล่นแต่ละคน
ทำให้ตรรกะของเกมชัดเจนและป้องกันการกระทำที่ไม่ถูกจังหวะ เช่น การร่ายคาถาซ้อนกัน

\section{การประมวลผลหลายเธรดใน Python}
Python มีกลไก Global Interpreter Lock (GIL) ที่อนุญาตให้รันโค้ด Python ได้ทีละเธรด
อย่างไรก็ตาม ไลบรารีที่เขียนด้วยภาษา C/C++ เช่น OpenCV และ MediaPipe จะปล่อย GIL ระหว่างการคำนวณหนัก
การแยกงานเหล่านี้ไปไว้คนละเธรดจึงทำงานขนานกันได้จริงบนหน่วยประมวลผลหลายแกน

% ==============================================================================
\chapter{การออกแบบระบบ}
% ==============================================================================

\section{ภาพรวมของเกม}
ผู้เล่นสองคนยืนหน้ากล้อง ผู้เล่นทางซ้ายของภาพเป็นผู้เล่น 1 (P1) และทางขวาเป็นผู้เล่น 2 (P2)
แต่ละคนมีพลังชีวิต (HP) 100 หน่วย และมานา (Mana) 100 หน่วย การร่ายคาถาหนึ่งครั้งมี 3 ขั้นตอน
\begin{enumerate}[leftmargin=2cm]
  \item \textbf{วาด} ชูนิ้วชี้นิ้วเดียวแล้ววาดสัญลักษณ์ของคาถากลางอากาศ จากนั้นหยุดนิ้วนิ่งเพื่อยืนยัน
  \item \textbf{ทำท่า} ทำท่ามือประจำคาถาเพื่อเริ่มชาร์จพลัง ระบบจะหักมานา ณ จุดนี้
  \item \textbf{เล็งและปล่อย} ระหว่างชาร์จ ทิศที่มือชี้คือทิศที่คาถาจะพุ่งไป เมื่อชาร์จครบเวลาคาถาจะถูกปล่อยอัตโนมัติ
\end{enumerate}
ผู้ที่ทำให้พลังชีวิตของอีกฝ่ายหมดก่อนเป็นผู้ชนะยกนั้น และผู้ที่ชนะ 2 ยกก่อนเป็นผู้ชนะการแข่งขัน

\begin{figure}[H]
  \centering
  \IfFileExists{screenshot_battle.png}{%
    \includegraphics[width=0.9\textwidth]{screenshot_battle.png}%
  }{%
    \fbox{\parbox[c][6cm][c]{0.85\textwidth}{\centering
      [ใส่ภาพหน้าจอขณะเล่นเกม]\\[0.2cm]
      บันทึกไฟล์ชื่อ \code{screenshot_battle.png} ไว้ในโฟลเดอร์เดียวกับไฟล์รายงานนี้}}%
  }
  \caption{หน้าจอขณะต่อสู้ แสดงแถบพลังชีวิต มานา จำนวนยกที่ชนะ และการ์ดคาถาของผู้เล่นแต่ละฝั่ง}
  \label{fig:screenshot}
\end{figure}

\section{คาถาและสมดุลของเกม}
ตารางที่~\ref{tab:spells} แสดงคุณสมบัติของคาถาทั้งหมด ซึ่งกำหนดไว้ในคลาส \code{SpellConfig} ของไฟล์ \code{spell_system.py}
การออกแบบให้คาถาที่รุนแรงกว่ามีต้นทุนมานา เวลาชาร์จ และคูลดาวน์สูงกว่า เพื่อให้ผู้เล่นต้องตัดสินใจเชิงกลยุทธ์

\begin{table}[H]
  \centering
  \caption{คุณสมบัติของคาถาในเกม}
  \label{tab:spells}
  \small
  \begin{tabularx}{\textwidth}{l l l r r r r Y}
    \toprule
    \textbf{คาถา} & \textbf{สัญลักษณ์} & \textbf{ท่าปล่อย} & \textbf{ดาเมจ} & \textbf{มานา} & \textbf{ชาร์จ (s)} & \textbf{คูลดาวน์ (s)} & \textbf{คุณสมบัติพิเศษ} \\
    \midrule
    Fireball         & สามเหลี่ยม & แบมือ      & 20  & 30 & 1.3  & 10 & -- \\
    Magic Shield     & วงกลม     & กำหมัด     & 0   & 25 & 0.8  & 15 & กางโล่ 5.5 วินาที \\
    Lightning Strike & ตัว Z      & ชูสองนิ้ว   & 40  & 50 & 2.5  & 30 & ความเร็วสูงสุด \\
    Terra Boulder    & สี่เหลี่ยม  & ชูสามนิ้ว   & 20  & 30 & 1.8  & 15 & ทะลุและทำลายโล่ \\
    Healing Light    & หัวใจ      & ชูนิ้วโป้ง  & 0   & 35 & 1.5  & 20 & ฟื้นพลังชีวิต 25 หน่วย \\
    Hollow Purple    & อนันต์     & ท่า OK     & 999 & 65 & 10.0 & 12 & คาถาพิเศษ ต้องปลดล็อก \\
    \bottomrule
  \end{tabularx}
\end{table}

\noindent กติกาเพิ่มเติมที่เกี่ยวกับทรัพยากรของผู้เล่นมีดังนี้
\begin{itemize}[leftmargin=2cm]
  \item มานาฟื้นฟูเอง 1.5 หน่วยต่อวินาที และ 8 หน่วยต่อวินาทีเมื่อทำท่าชาร์จมานา (ชูสี่นิ้วโดยพับนิ้วโป้ง)
  \item มานาไม่ฟื้นฟูระหว่างกางโล่หรือระหว่างชาร์จคาถา
  \item หลังวาดสัญลักษณ์สำเร็จ ผู้เล่นมีเวลา 4 วินาทีในการทำท่าปล่อย มิฉะนั้นคาถาจะถูกยกเลิก
        และยกเลิกเองได้ด้วยการกำหมัดหรือลดมือลง
\end{itemize}

\subsection*{คาถาพิเศษ Hollow Purple}
คาถานี้ได้แรงบันดาลใจจากการ์ตูนเรื่อง Jujutsu Kaisen ออกแบบเป็นกลไกพลิกเกมสำหรับฝ่ายที่กำลังเสียเปรียบ
ผู้เล่นจะปลดล็อกได้เมื่อครบเงื่อนไขทั้งหมดต่อไปนี้ (1) เคยร่ายคาถาพื้นฐานครบทั้ง 5 ชนิดในยกนั้น
(2) พลังชีวิตต่ำกว่า 50 (3) มานาไม่น้อยกว่า 40 และ (4) ทำท่ามือต่อเนื่อง กำหมัด $\rightarrow$ ชูสองนิ้ว $\rightarrow$ แบมือ
โดยแต่ละท่าห่างกันไม่เกิน 4 วินาที เมื่อปลดล็อกแล้วผู้เล่นยังใช้บาเรีย ``Infinity'' ได้ด้วยการชูนิ้วชี้และนิ้วกลางชิดกัน
ซึ่งป้องกันคาถาทุกชนิดยกเว้น Hollow Purple

\section{สถาปัตยกรรมของระบบ}
ระบบแบ่งการทำงานเป็น 3 ส่วนที่ทำงานขนานกัน ดังรูปที่~\ref{fig:arch} ได้แก่ เธรดอ่านภาพจากกล้อง
เธรดประมวลผลโมเดล และลูปหลักของเกม ทั้งสามส่วนสื่อสารกันโดยส่งต่อเฉพาะ ``ข้อมูลล่าสุด'' ไม่มีการเข้าคิว
จึงไม่เกิดความล่าช้าสะสม

\begin{figure}[H]
  \centering
  \begin{tikzpicture}[node distance=0.9cm and 1.1cm, font=\small]
    \node[box] (cam) {กล้องเว็บแคม};
    \node[thr, right=of cam] (stream) {CameraStream\\(เธรดอ่านภาพ)};
    \node[thr, right=2.6cm of stream] (vision) {VisionWorker\\(เธรดโมเดล)};
    \node[thr, below=1.1cm of vision, xshift=-1.6cm] (p1) {PlayerVision P1\\มือ + ใบหน้า};
    \node[thr, below=1.1cm of vision, xshift=1.9cm] (p2) {PlayerVision P2\\มือ + ใบหน้า};
    \node[box, below=3.4cm of stream] (logic) {ลูปหลัก: ตรรกะเกม\\ท่ามือ / ลายเส้น / การชน};
    \node[box, left=of logic] (render) {UIRenderer\\เอฟเฟกต์และ HUD};
    \node[box, below=of render] (screen) {หน้าจอ};
    \node[box, right=of logic] (sound) {SoundManager\\เสียง};
    \draw[arr] (cam) -- (stream);
    \draw[arr] (stream) -- node[above] {เฟรมล่าสุด} (vision);
    \draw[arr] (vision) -- (p1);
    \draw[arr] (vision) -- (p2);
    \draw[arr] (stream) -- node[left, align=right] {เฟรม\\ล่าสุด} (logic);
    \draw[arr] (p1) -- node[right, pos=0.45] {\ พิกัดมือและใบหน้า} (logic);
    \draw[arr] (logic) -- (render);
    \draw[arr] (render) -- (screen);
    \draw[arr] (logic) -- (sound);
  \end{tikzpicture}
  \caption{สถาปัตยกรรมของระบบ กล่องสีฟ้าคือส่วนที่ทำงานบนเธรดแยกจากลูปหลัก}
  \label{fig:arch}
\end{figure}

\section{โครงสร้างไฟล์ของโปรแกรม}
\begin{table}[H]
  \centering
  \caption{ไฟล์ซอร์สโค้ดและหน้าที่}
  \label{tab:files}
  \small
  \begin{tabularx}{\textwidth}{l Y}
    \toprule
    \textbf{ไฟล์} & \textbf{หน้าที่} \\
    \midrule
    \code{game_engine.py} & ลูปหลักของเกม ประสานทุกระบบ เครื่องสถานะการแข่งขัน ฟิสิกส์ของกระสุนและการตรวจการชน แถบสถานะบนจอ \\
    \code{vision_pipeline.py} & เธรดอ่านภาพจากกล้อง และเธรดประมวลผล MediaPipe แยกตามผู้เล่น \\
    \code{gesture_recognizer.py} & จำแนกท่ามือจากจุดสำคัญ 21 จุด และคำนวณทิศทางการเล็ง \\
    \code{unistroke.py} & รู้จำลายเส้นที่วาดกลางอากาศด้วย \$1 Recognizer \\
    \code{spell_system.py} & ค่าสมดุลของคาถา และเครื่องสถานะของผู้เล่น (มานา คูลดาวน์ โล่) \\
    \code{ui_renderer.py} & ระบบอนุภาค การเบลนด์ภาพเอฟเฟกต์ แอนิเมชันชาร์จคาถา การ์ดคาถา \\
    \code{sound_manager.py} & เพลงประกอบ เสียงเอฟเฟกต์ เสียงวนต่อเนื่อง และเสียงผู้ประกาศ \\
    \code{config.py} & ค่าตั้งทั้งหมดของเกมในที่เดียว (กล้อง โมเดล กติกา จังหวะเวลา) \\
    \code{resource_helper.py} & แปลงเส้นทางไฟล์ทรัพยากร รองรับทั้งการรันด้วย Python และไฟล์ .exe \\
    \code{test_*.py} & สคริปต์ทดสอบแยกส่วนตามระยะการพัฒนา (ตรวจจับมือ การเล็ง การวาด) \\
    \code{fix_audio_pcm.py} & เครื่องมือแปลงไฟล์เสียง WAV เป็นรูปแบบ 16-bit PCM \\
    \bottomrule
  \end{tabularx}
\end{table}

\section{เครื่องสถานะของการแข่งขัน}
การแข่งขันมี 5 สถานะ ดังรูปที่~\ref{fig:matchfsm} ที่หน้าไตเติล ผู้เล่นทั้งสองต้องชูนิ้วโป้งหรือชูสองนิ้วค้างพร้อมกัน 1.2 วินาที
เพื่อเริ่มเกม ซึ่งเป็นการยืนยันว่าระบบมองเห็นมือของทั้งสองฝั่งแล้วจริง หลังจบการแข่งขันใช้การชูนิ้วโป้งค้าง 1.5 วินาทีเพื่อแข่งใหม่

\begin{figure}[H]
  \centering
  \begin{tikzpicture}[font=\small]
    \node[state] (title) at (0,0) {TITLE};
    \node[state] (count) at (4.6,0) {COUNTDOWN};
    \node[state] (battle) at (9.2,0) {BATTLE};
    \node[state] (round) at (6.9,-2.6) {ROUND\_OVER};
    \node[state] (match) at (11.7,-2.6) {MATCH\_OVER};
    \draw[arr] (title) -- node[above] {\footnotesize ทั้งคู่พร้อม} (count);
    \draw[arr] (count) -- node[above] {\footnotesize นับถอยหลังจบ} (battle);
    \draw[arr] (battle) -- node[right, pos=0.5] {\ HP หมด} (round);
    \draw[arr] (round) -| node[right, pos=0.72] {\footnotesize พัก 2.4 วินาที} (count);
    \draw[arr] (battle) -| node[right, pos=0.75, align=left] {ชนะครบ\\2 ยก} (match);
    \draw[arr] (match) -- ++(0,-1.4) -| node[below, pos=0.25] {ทั้งคู่ขอแข่งใหม่} (count.south west);
  \end{tikzpicture}
  \caption{เครื่องสถานะของการแข่งขัน}
  \label{fig:matchfsm}
\end{figure}

\section{เครื่องสถานะของผู้เล่น}
คลาส \code{PlayerCombatState} จัดการสถานะการร่ายคาถาของผู้เล่นแต่ละคนแยกจากกัน ดังรูปที่~\ref{fig:playerfsm}

\begin{figure}[H]
  \centering
  \begin{tikzpicture}[font=\small]
    \node[state] (idle) at (0,0) {IDLE};
    \node[state] (rec) at (5.5,0) {RECOGNIZED};
    \node[state] (chg) at (11,0) {CHARGING};
    \draw[arr] (idle) -- node[below, align=center] {วาดสัญลักษณ์ถูกต้อง\\มานาพอ ไม่ติดคูลดาวน์} (rec);
    \draw[arr] (rec) -- node[above, align=center] {ทำท่าปล่อย\\(หักมานา)} (chg);
    \draw[arr] (chg) -- ++(0,-2.3) -| node[below, pos=0.25] {ชาร์จครบเวลา: ปล่อยคาถา เริ่มคูลดาวน์} (idle);
    \draw[arr] (rec) -- ++(0,1.5) -| node[above, pos=0.25] {หมดเวลา 4 วินาที / กำหมัด / ลดมือ} (idle);
  \end{tikzpicture}
  \caption{เครื่องสถานะการร่ายคาถาของผู้เล่น}
  \label{fig:playerfsm}
\end{figure}

% ==============================================================================
\chapter{การพัฒนาระบบ}
% ==============================================================================

\section{การรับภาพและการแบ่งแดนผู้เล่น}
ภาพจากกล้องถูกอ่านโดยคลาส \code{CameraStream} ซึ่งเก็บเฉพาะเฟรมล่าสุด ผลการตรวจจับมือแต่ละมือจะถูกจัดให้ผู้เล่น
ตามตำแหน่งข้อมือในแนวนอน หากข้อมืออยู่ครึ่งซ้ายของภาพถือเป็นของ P1 มิฉะนั้นเป็นของ P2
เมื่อผู้เล่นคนหนึ่งมีมือมากกว่าหนึ่งมือในภาพ ระบบเลือก ``มือที่ยื่นเข้าหาคู่ต่อสู้มากที่สุด''
(มือที่อยู่ใกล้เส้นกึ่งกลางจอที่สุด) เป็นมือที่ใช้ควบคุม ทำให้มืออีกข้างที่ปล่อยไว้ข้างลำตัวไม่รบกวนการเล่น

\section{การจำแนกท่ามือ}
คลาส \code{GestureRecognizer} จำแนกท่ามือด้วยกฎเชิงเรขาคณิต ไม่ได้ฝึกโมเดลเพิ่ม ขั้นแรกคำนวณ \textit{ขนาดฝ่ามือ}
จากระยะระหว่างข้อมือกับโคนนิ้วกลาง แล้วใช้เป็นหน่วยอ้างอิงของระยะทั้งหมด ทำให้เกณฑ์ไม่ขึ้นกับระยะห่างจากกล้อง
\begin{equation}
  p = \max\bigl(0.015,\; \lVert \mathbf{w} - \mathbf{m} \rVert \bigr)
\end{equation}
โดย $\mathbf{w}$ และ $\mathbf{m}$ คือพิกัดของข้อมือและโคนนิ้วกลาง จากนั้นตัดสินว่านิ้วแต่ละนิ้ว ``เหยียด'' หรือ ``งอ''
โดยนิ้วชี้ถึงนิ้วก้อยถือว่าเหยียดเมื่อปลายนิ้วอยู่สูงกว่าข้อกลางนิ้ว หรือปลายนิ้วอยู่ไกลจากข้อมือมากกว่าข้อกลางนิ้วและโคนนิ้วตามสัดส่วนที่กำหนด
ดังโค้ดต่อไปนี้

\begin{lstlisting}[caption={การตัดสินสถานะนิ้วชี้ถึงนิ้วก้อย (ย่อจาก gesture\_recognizer.py)}, label={lst:finger}]
for tip, pip, mcp in zip(tips, pips, mcps):
    d_wrist_tip = math.hypot(lm[tip].x - wrist.x, lm[tip].y - wrist.y)
    d_wrist_pip = math.hypot(lm[pip].x - wrist.x, lm[pip].y - wrist.y)
    d_wrist_mcp = math.hypot(lm[mcp].x - wrist.x, lm[mcp].y - wrist.y)

    is_pointing_up = (lm[tip].y < lm[pip].y)
    is_extended_out = (d_wrist_tip > d_wrist_pip * 1.06) and \
                      (d_wrist_tip > d_wrist_mcp * 1.20)
    fingers.append(is_pointing_up or is_extended_out)
\end{lstlisting}

เมื่อได้สถานะของนิ้วทั้งห้าแล้วจึงจับคู่กับท่ามือตามตารางที่~\ref{tab:gestures} ท่า OK และท่าชูนิ้วโป้งมีการตรวจเฉพาะ
และตรวจก่อนท่าอื่น เพราะรูปร่างของสองท่านี้อาจถูกตีความเป็นท่าอื่นได้หากดูเพียงจำนวนนิ้วที่เหยียด

\begin{table}[H]
  \centering
  \caption{ท่ามือที่ระบบจำแนกและการใช้งาน}
  \label{tab:gestures}
  \small
  \begin{tabularx}{\textwidth}{l Y Y}
    \toprule
    \textbf{ท่ามือ} & \textbf{เงื่อนไขโดยสรุป} & \textbf{การใช้งานในเกม} \\
    \midrule
    AIR\_DRAW & เหยียดนิ้วชี้นิ้วเดียว & วาดสัญลักษณ์ \\
    FIST & งอนิ้วชี้ถึงนิ้วก้อยทั้งหมด & ปล่อย Magic Shield / ยกเลิกคาถา \\
    PEACE & เหยียดนิ้วชี้และนิ้วกลาง & ปล่อย Lightning Strike / พร้อมเริ่มเกม \\
    THREE\_FINGERS & เหยียดนิ้วชี้ กลาง นาง & ปล่อย Terra Boulder \\
    OPEN\_HAND & เหยียดครบห้านิ้ว & ปล่อย Fireball \\
    RECHARGE & เหยียดสี่นิ้ว พับนิ้วโป้ง & ชาร์จมานา \\
    THUMBS\_UP & นิ้วโป้งชี้ขึ้นเหนือข้อนิ้วอื่น และนิ้วอื่นงออย่างน้อย 3 นิ้ว & ปล่อย Healing Light / พร้อมเริ่มเกม / ขอแข่งใหม่ \\
    OK\_SIGN & ปลายนิ้วโป้งแตะปลายนิ้วชี้ และนิ้วที่เหลือเหยียด & ปล่อย Hollow Purple \\
    \bottomrule
  \end{tabularx}
\end{table}

\section{การวาดเส้นกลางอากาศ}
คลาส \code{AirDrawSession} เก็บเส้นทางของปลายนิ้วชี้ขณะผู้เล่นทำท่า AIR\_DRAW โดยมีรายละเอียดดังนี้
\begin{itemize}[leftmargin=2cm]
  \item \textbf{การกรองการสั่น} พิกัดปลายนิ้วผ่านการกรองตามสมการ~\eqref{eq:ema} ด้วย $\alpha = 0.75$
        และบันทึกจุดใหม่เฉพาะเมื่อห่างจากจุดก่อนหน้าอย่างน้อย 4 พิกเซล
  \item \textbf{การยืนยันด้วยการหยุดนิ่ง} ระบบประมาณความเร็วของปลายนิ้ว เมื่อความเร็วต่ำกว่าเกณฑ์ต่อเนื่อง 0.55 วินาที
        ถือว่าวาดเสร็จ ระหว่างรอจะแสดงวงแหวนความคืบหน้าที่ปลายนิ้วเพื่อให้ผู้เล่นรู้ว่าระบบกำลังยืนยัน
  \item \textbf{ความทนทานต่อการตรวจจับผิดชั่วขณะ} หากท่ามือหลุดจาก AIR\_DRAW ชั่วครู่ ระบบจะยังไม่ตัดเส้นทันที
        แต่รอช่วงเวลาผ่อนผันก่อน เพื่อไม่ให้เส้นขาดเมื่อวาดเร็ว
  \item \textbf{การป้องกันการวาดซ้ำโดยไม่ตั้งใจ} หลังวาดเสร็จผู้เล่นต้องเปลี่ยนเป็นท่าอื่นก่อนจึงจะเริ่มเส้นใหม่ได้
\end{itemize}

\section{การรู้จำลายเส้น}
คลาส \code{UnistrokeRecognizer} พัฒนาตามหลักการของ \$1 Recognizer โดยใช้ $N = 64$ จุด และกรอบมาตรฐานขนาด 250 หน่วย
มีแม่แบบของสัญลักษณ์ 6 ชนิด บางชนิดมีมากกว่าหนึ่งแม่แบบเพื่อรองรับลักษณะการวาดที่ต่างกัน
(เช่น สี่เหลี่ยมมุมแหลมและมุมมน) ผู้จัดทำได้ดัดแปลงจากอัลกอริทึมต้นฉบับดังนี้

\begin{enumerate}[leftmargin=2cm]
  \item \textbf{รองรับการวาดย้อนทิศ} เปรียบเทียบทั้งลายเส้นตามลำดับที่วาดและลำดับย้อนกลับ ผู้เล่นจึงวาดตามหรือทวนเข็มนาฬิกาก็ได้
  \item \textbf{รองรับจุดเริ่มต้นใดก็ได้} สำหรับรูปปิด ทดสอบการเลื่อนจุดเริ่มต้นแบบวนรอบ (Cyclic Shift) 32 ตำแหน่ง
        ผู้เล่นจึงเริ่มวาดวงกลมหรือสามเหลี่ยมจากจุดใดก็ได้
  \item \textbf{จำกัดมุมหมุน} ค้นหามุมที่ดีที่สุดในช่วง $\pm 45^\circ$ เท่านั้น เพื่อไม่ให้รูปที่ต่างกันเพียงการหมุนถูกตีความปนกัน
  \item \textbf{กรองแม่แบบด้วยลักษณะทางเรขาคณิต} ก่อนเปรียบเทียบจะคำนวณ (ก) ลายเส้นเป็นรูปปิดหรือไม่
        จากระยะระหว่างจุดเริ่มและจุดจบเทียบกับเส้นทแยงมุมของกรอบ (ข) พื้นที่ของรูปด้วยสูตร Shoelace
        และ (ค) จำนวนมุมหักจากผลคูณจุดของเวกเตอร์ แล้วตัดแม่แบบที่เป็นไปไม่ได้ออก
        เช่น รูปที่มีพื้นที่เกินร้อยละ 60 ของกรอบจะไม่ถูกเทียบกับสามเหลี่ยม วิธีนี้ช่วยแยกวงกลมออกจากสี่เหลี่ยมมุมมนซึ่งมีลายเส้นคล้ายกันมาก
\end{enumerate}

ผลลัพธ์จะถูกยอมรับเมื่อคะแนนตามสมการ~\eqref{eq:score} ไม่น้อยกว่า 0.65 หากต่ำกว่านั้นระบบแสดงข้อความ ``Uncertain''
ให้ผู้เล่นวาดใหม่

\section{การเล็งและกระสุนเวทมนตร์}
ทิศทางการเล็งคำนวณจากเวกเตอร์ที่ลากจากข้อมือไปยังปลายนิ้วกลาง โดยบังคับให้องค์ประกอบแนวนอนชี้เข้าหาคู่ต่อสู้เสมอ
(ไปทางขวาสำหรับ P1 และทางซ้ายสำหรับ P2) ผู้เล่นจึงปรับได้เฉพาะมุมเงยและมุมก้ม ซึ่งป้องกันการยิงผิดฝั่งโดยไม่ตั้งใจ
ระหว่างชาร์จจะมีเส้นนำสายตาแสดงทิศทางปัจจุบัน

เมื่อชาร์จครบ ระบบสร้างวัตถุ \code{Projectile} ซึ่งเคลื่อนที่ด้วยความเร็วคงที่ตามเวลาจริงที่ผ่านไปในแต่ละเฟรม
\begin{equation}
  \mathbf{x}_{t+\Delta t} = \mathbf{x}_t + \mathbf{v}\,\Delta t
\end{equation}
การคูณด้วย $\Delta t$ ทำให้ความเร็วของกระสุนบนจอเท่าเดิมไม่ว่าเกมจะแสดงผลได้กี่เฟรมต่อวินาที

\section{การตรวจการชนและเป้ารับความเสียหาย}
เป้ารับความเสียหายของผู้เล่นเป็นกรอบสี่เหลี่ยมบริเวณลำตัว ซึ่งคลาส \code{TorsoTracker} ประมาณจากตำแหน่งใบหน้า
โดยเลื่อนลงมาตามระยะที่กำหนดและกรองด้วยสมการ~\eqref{eq:ema} เมื่อมองไม่เห็นใบหน้าชั่วคราว (เช่น ผู้เล่นก้มหลบ)
ระบบจะคงตำแหน่งเดิมไว้ระยะหนึ่ง การตรวจการชนในแต่ละเฟรมมีลำดับดังนี้
\begin{enumerate}[leftmargin=2cm]
  \item \textbf{ชนโล่} หากฝ่ายรับกางโล่อยู่ และระยะจากกระสุนถึงศูนย์กลางโล่ไม่เกินรัศมีโล่บวกรัศมีกระสุน
        กระสุนทั่วไปจะถูกป้องกัน ส่วนกระสุนที่มีคุณสมบัติทะลุโล่จะทำลายโล่และสร้างความเสียหาย
  \item \textbf{ชนลำตัว} หาจุดบนกรอบลำตัวที่ใกล้กระสุนที่สุด หากระยะไม่เกินรัศมีกระสุนบวกระยะเผื่อ ถือว่าโดน
  \item \textbf{บาเรีย Infinity} หากฝ่ายรับใช้บาเรียอยู่ กระสุนทุกชนิดยกเว้น Hollow Purple จะถูกลบล้าง
\end{enumerate}
การให้ผู้เล่น ``หลบ'' ได้จริงด้วยการขยับตัวออกจากแนวกระสุนเป็นจุดเด่นของการใช้ตำแหน่งร่างกายเป็นเป้า

\section{ระบบแสดงผลและเอฟเฟกต์}
คลาส \code{UIRenderer} รับผิดชอบการวาดทั้งหมดลงบนเฟรมภาพจากกล้อง ประกอบด้วย
\begin{itemize}[leftmargin=2cm]
  \item \textbf{ระบบอนุภาค} มีอนุภาคหลายชนิด เช่น ประกายไฟ เศษหิน ประกายไฟฟ้า ควัน และหัวใจลอยขึ้น
        แต่ละชนิดมีตำแหน่ง ความเร็ว และอายุของตนเอง และอัปเดตตามเวลาจริงที่ผ่านไป
  \item \textbf{แอนิเมชันชาร์จคาถา} คาถาแต่ละชนิดมีเอฟเฟกต์ชาร์จเฉพาะตัวที่ประกอบจากภาพหลายชั้นซึ่งหมุนและขยายตามความคืบหน้า
        พร้อมป้ายนับเวลาถอยหลัง
  \item \textbf{เอฟเฟกต์การปะทะ} แรงระเบิด คลื่นกระแทก และการสั่นของหน้าจอซึ่งมีความแรงต่างกันตามชนิดของคาถา
  \item \textbf{ฉากปิดยก} เมื่อพลังชีวิตของฝ่ายใดหมด เกมจะเล่นแบบสโลว์โมชัน (ร้อยละ 25 ของความเร็วปกติ) เป็นเวลา 0.75 วินาที
        พร้อมแสงวาบสีขาว
  \item \textbf{ส่วนแสดงสถานะ (HUD)} แถบพลังชีวิต แถบมานา จำนวนยกที่ชนะ และการ์ดคาถาที่แสดงสัญลักษณ์ ท่าปล่อย ต้นทุนมานา
        และเวลาคูลดาวน์ที่เหลือ
\end{itemize}

\section{ระบบเสียง}
คลาส \code{SoundManager} ใช้ Pygame mixer โดยจัดสรรช่องเสียง 32 ช่อง และ \textit{จองช่อง 8 ช่องแรก} ไว้สำหรับเสียงที่เล่นวนต่อเนื่อง
ของผู้เล่นแต่ละคน ได้แก่ เสียงวาด เสียงชาร์จ เสียงโล่ และเสียงชาร์จมานา การจองช่องแยกเช่นนี้ป้องกันปัญหาเสียงวนถูกเสียงเอฟเฟกต์อื่นแย่งช่อง
แล้วขาดหาย ส่วนเสียงผู้ประกาศ (เช่น ``Round 1'' และ ``K.O.'') ใช้คิวที่ตั้งเวลาหน่วงได้ เพื่อให้จังหวะเสียงตรงกับภาพ

\section{การรวมค่าตั้งไว้ที่เดียว}
ค่าคงที่ทั้งหมดของเกม เช่น ความละเอียดกล้อง ค่าความเชื่อมั่นของโมเดล ขนาดเป้ารับความเสียหาย จังหวะเวลา และความแรงของการสั่นหน้าจอ
รวมไว้ในไฟล์ \code{config.py} ทำให้ปรับสมดุลของเกมและปรับให้เข้ากับฮาร์ดแวร์ได้โดยไม่ต้องแก้ตรรกะของโปรแกรม

% ==============================================================================
\chapter{การปรับปรุงประสิทธิภาพ}
% ==============================================================================

เกมที่ควบคุมด้วยการเคลื่อนไหวต้องตอบสนองทันที ผู้จัดทำจึงวัดประสิทธิภาพของโปรแกรมรุ่นแรกและพบว่าแสดงผลได้เพียง
17--30 เฟรมต่อวินาที บทนี้อธิบายวิธีวัด สาเหตุที่พบ และการแก้ไข

\section{วิธีการวัด}
ผู้จัดทำเขียนสคริปต์ที่รันเกมจริงโดยปิดเสียง จำลองการกดปุ่มเพื่อเข้าสู่แต่ละสถานะ และบันทึกเวลาของทุกเฟรม
สำหรับกรณีภาระสูงสุด สคริปต์จะสั่งให้ผู้เล่นทั้งสองฝั่งร่ายคาถาต่อเนื่องทุก 1.6 วินาทีโดยวนครบทุกชนิด
การวัดทำบนคอมพิวเตอร์ที่ใช้พัฒนา (หน่วยประมวลผล 12 เธรด ระบบปฏิบัติการ Windows 10 ไม่ใช้หน่วยประมวลผลกราฟิกในการคำนวณ)
ตัวเลขจึงเป็นค่าเฉพาะของเครื่องนี้และใช้เพื่อเปรียบเทียบก่อนและหลังการแก้ไขเป็นหลัก

\section{คอขวดที่พบและการแก้ไข}

\subsection{กล้องถูกจำกัดที่ 30 เฟรมต่อวินาที}
ที่ความละเอียด 1280$\times$720 กล้องส่งภาพแบบไม่บีบอัดได้เพียง 30 เฟรมต่อวินาที เมื่อขอรูปแบบ MJPG
(ภาพบีบอัดแบบ JPEG ทีละเฟรม) กล้องตัวเดียวกันส่งได้ประมาณ 60 เฟรมต่อวินาที

\subsection{การอ่านกล้องและโมเดลอยู่ในลูปเดียวกับการวาด}
ในรุ่นแรก ทุกรอบของลูปต้องรอภาพจากกล้อง แล้วรอโมเดลตรวจจับมือ (ประมาณ 20 มิลลิวินาที) และโมเดลตรวจจับใบหน้า (ประมาณ 8 มิลลิวินาที)
ก่อนจึงจะวาดภาพได้ อัตราการแสดงผลจึงถูกจำกัดด้วยผลรวมของเวลาทั้งหมด การแก้ไขคือแยกเป็นสามส่วนตามรูปที่~\ref{fig:arch}
ลูปหลักจะรอเฟรมใหม่จากกล้องไม่เกินเวลาที่กำหนด หากยังไม่มีเฟรมใหม่ก็วาดเอฟเฟกต์ซ้ำบนเฟรมล่าสุด
เอฟเฟกต์และแถบสถานะจึงเคลื่อนไหวลื่นไหลโดยไม่ขึ้นกับความเร็วของโมเดล

การแยกเธรดทำให้ต้องแก้ตรรกะที่เคยนับ ``ต่อเฟรม'' สองจุด คือ (1) การประมาณความเร็วปลายนิ้วในการวาด
ซึ่งต้องคำนวณเฉพาะเมื่อมีผลการตรวจจับชุดใหม่และปรับตามช่วงเวลาจริง มิฉะนั้นตำแหน่งซ้ำจะถูกตีความว่านิ้วหยุดนิ่ง
และ (2) อัตราการปล่อยอนุภาค ซึ่งกำหนดให้คงที่ 30 ครั้งต่อวินาที เพื่อไม่ให้เอฟเฟกต์หนาแน่นขึ้นตามอัตราการแสดงผล

\subsection{การเบลนด์ภาพเอฟเฟกต์}
ภาพเอฟเฟกต์แต่ละชิ้นถูกย่อขยาย หมุน แยกช่องความทึบ คูณ แล้วจึงบวกเข้ากับเฟรม รวมหลายขั้นตอนต่อชิ้นต่อเฟรม
ผู้จัดทำแก้ไขโดยคำนวณภาพที่คูณค่าความทึบไว้ล่วงหน้า (Premultiplied Alpha) เพียงครั้งเดียวแล้วเก็บไว้ใช้ซ้ำ
รวมการย่อขยายกับการหมุนเป็นการแปลงครั้งเดียวเมื่อเป็นการขยายภาพ และประมวลผลเฉพาะบริเวณที่มองเห็นบนจอ
เวลาในการวาดเอฟเฟกต์ที่หนักที่สุด (Hollow Purple ของผู้เล่นสองฝั่งพร้อมกัน) ลดจากประมาณ 78 เหลือประมาณ 28 มิลลิวินาทีต่อเฟรม
ผู้จัดทำเปรียบเทียบภาพผลลัพธ์ก่อนและหลังแก้ไขแบบพิกเซลต่อพิกเซล พบความต่างเฉพาะที่ขอบของลายเส้นบาง ๆ ซึ่งมองด้วยตาไม่เห็น

\subsection{ตัวรู้จำลายเส้นทำให้เกมหยุดชะงัก}
การรู้จำหนึ่งครั้งต้องเปรียบเทียบกับแม่แบบหลายอัน แต่ละอันทดสอบ 32 จุดเริ่มต้น $\times$ 2 ทิศทาง
และแต่ละครั้งต้องค้นหามุมที่ดีที่สุดอีกประมาณ 10 รอบ เมื่อเขียนด้วยลูปของ Python เกมจึงหยุดนิ่งประมาณ 270 มิลลิวินาทีทุกครั้งที่ผู้เล่นวาดเสร็จ
ผู้จัดทำแก้ไขโดยใช้ NumPy คำนวณทุกจุดเริ่มต้นและทั้งสองทิศทางพร้อมกันเป็นชุดเดียว ดังโค้ดที่~\ref{lst:batch}
โดยคงขั้นตอนของอัลกอริทึมไว้เหมือนเดิมทุกประการ

\begin{lstlisting}[caption={ฟังก์ชันระยะห่างที่คำนวณลายเส้นหลายชุดพร้อมกัน (ย่อจาก unistroke.py)}, label={lst:batch}]
# dx, dy: (B, N) offsets from each candidate's centroid
# angles: (B,) one trial rotation per candidate
def distance(angles):
    cos_val = np.cos(angles)[:, np.newaxis]
    sin_val = np.sin(angles)[:, np.newaxis]
    rx = dx * cos_val - dy * sin_val + cx
    ry = dx * sin_val + dy * cos_val + cy
    return np.hypot(rx - tx, ry - ty).mean(axis=1)
\end{lstlisting}

เมื่อทดสอบกับลายเส้นสังเคราะห์ 400 เส้น (หลายรูปทรง หมุน ย่อขยาย และเติมสัญญาณรบกวนแบบสุ่ม)
ผลการจำแนกของรุ่นใหม่ตรงกับรุ่นเดิมทุกเส้น และคะแนนต่างกันไม่เกิน $10^{-15}$

\section{ผลการปรับปรุง}
\begin{table}[H]
  \centering
  \caption{อัตราการแสดงผลก่อนและหลังการปรับปรุง (เฟรมต่อวินาที)}
  \label{tab:fps}
  \begin{tabularx}{\textwidth}{Y r r}
    \toprule
    \textbf{สถานการณ์} & \textbf{ก่อน} & \textbf{หลัง} \\
    \midrule
    หน้าไตเติล & 16.7 & 87 \\
    นับถอยหลัง & 30.0 & 86 \\
    ต่อสู้ ภาระสูงสุด (สองฝั่งร่ายคาถาต่อเนื่อง) & 22.9 & 65 \\
    ต่อสู้ ไม่มีเอฟเฟกต์บนจอ & ไม่ได้วัด & 88 \\
    \midrule
    อัตราเฟรมของกล้อง & 30 & ประมาณ 60 \\
    เวลาที่เกมหยุดเมื่อรู้จำลายเส้น (มิลลิวินาที) & ประมาณ 270 & ประมาณ 8 \\
    \bottomrule
  \end{tabularx}
\end{table}

ข้อสังเกตคือ อัตราการตรวจจับมือ (ประมาณ 25--30 ครั้งต่อวินาที) ยังต่ำกว่าอัตราการแสดงผล
เนื่องจากถูกจำกัดด้วยความเร็วของโมเดลบนหน่วยประมวลผลกลาง ตำแหน่งของเอฟเฟกต์ที่ติดกับมือจึงอาจตามหลังภาพจากกล้องเล็กน้อย

\section{การตรวจจับมือขนาดเล็กเมื่อผู้เล่นยืนไกล}
\label{sec:smallhand}
เมื่อผู้เล่นยืนไกลพอให้กล้องเห็นทั้งลำตัว มือในภาพจะเล็กจนระบบตรวจไม่พบ สาเหตุคือโมเดลตรวจจับฝ่ามือย่อภาพที่รับเข้าไปเป็นขนาดเล็กคงที่
เมื่อส่งทั้งเฟรมที่มีอัตราส่วน 16:9 เข้าไป มือจึงเหลือเพียงไม่กี่พิกเซล ผู้จัดทำแก้ไขดังนี้
\begin{enumerate}[leftmargin=2cm]
  \item \textbf{แยกตัวตรวจจับตามผู้เล่น} แต่ละฝั่งมีโมเดลของตนเองและทำงานขนานกัน โดยรับภาพที่ตัดเฉพาะแดนของผู้เล่น
        เป็นกรอบสี่เหลี่ยมจัตุรัส 720$\times$720 พิกเซลที่ความละเอียดจริง กรอบของสองฝั่งซ้อนกันบริเวณกลางจอ
        เพื่อไม่ให้มือที่ยื่นข้ามเส้นแบ่งแดนถูกตัดขาด
  \item \textbf{ซูมตามขนาดใบหน้า} ใช้ความสูงของใบหน้าเป็นตัวประมาณระยะของผู้เล่น แล้วกำหนดกรอบตรวจจับมือให้มีด้านยาว 10 เท่าของความสูงใบหน้า
        ซึ่งกว้างพอครอบคลุมระยะเอื้อมแขน ยิ่งผู้เล่นยืนไกล กรอบยิ่งแคบ มือจึงมีสัดส่วนในภาพที่ส่งให้โมเดลใกล้เคียงเดิม
  \item \textbf{ป้องกันกรอบสั่น} กรอบจะขยับเฉพาะเมื่อผู้เล่นเคลื่อนที่เกินร้อยละ 12 ของขนาดกรอบ
        เนื่องจากโมเดลติดตามมือโดยอาศัยตำแหน่งจากเฟรมก่อนหน้า กรอบที่ขยับทุกเฟรมจะทำให้การติดตามหลุดบ่อย
  \item \textbf{แปลงพิกัดกลับ} พิกัดที่ได้ถูกแปลงกลับเป็นพิกัดของเฟรมเต็ม ส่วนอื่นของโปรแกรมจึงไม่ต้องแก้ไข
\end{enumerate}

ผู้จัดทำทดสอบโดยนำภาพผู้เล่นที่ยกมือหนึ่งภาพมาย่อให้เล็กลงเป็นลำดับ วางในเฟรมขนาด 1280$\times$720 ทั้งสองฝั่ง แล้วส่งผ่านระบบตรวจจับจริง
ได้ผลดังตารางที่~\ref{tab:smallhand} ระบบเดิมตรวจไม่พบเมื่อมือสูงต่ำกว่าประมาณ 100 พิกเซล ส่วนระบบใหม่ตรวจพบได้ถึงประมาณ 50 พิกเซล
ข้อแลกเปลี่ยนคืออัตราการตรวจจับมือลดลงเล็กน้อย เนื่องจากต้องรันโมเดลสองชุด

\begin{table}[H]
  \centering
  \caption{ผลการตรวจจับมือที่ขนาดต่าง ๆ (ความสูงของมือในเฟรม 1280$\times$720)}
  \label{tab:smallhand}
  \begin{tabular}{r l l}
    \toprule
    \textbf{ความสูงของมือ (พิกเซล)} & \textbf{ระบบเดิม} & \textbf{ระบบใหม่} \\
    \midrule
    104 ขึ้นไป & พบ & พบ \\
    52--87 & ไม่พบ & พบทั้งสองฝั่ง \\
    43 & ไม่พบ & พบเพียงฝั่งเดียว \\
    37 & ไม่พบ & ไม่พบ \\
    \bottomrule
  \end{tabular}
\end{table}

% ==============================================================================
\chapter{การทดสอบและอภิปรายผล}
% ==============================================================================

\section{วิธีการทดสอบ}
\begin{itemize}[leftmargin=2cm]
  \item \textbf{การทดสอบแยกส่วนตามระยะการพัฒนา} ผู้จัดทำพัฒนาเป็นระยะและมีสคริปต์ทดสอบของแต่ละระยะ ได้แก่
        \code{test_hand_tracking.py} (การตรวจจับมือและการแบ่งแดน) \code{test_phase2_aiming.py} (การเลือกมือและการเล็ง)
        และ \code{test_phase3_air_drawing.py} (การวาดและการรู้จำลายเส้น) ซึ่งแสดงผลให้ตรวจด้วยตาผ่านกล้องจริง
  \item \textbf{การทดสอบวงจรการร่ายคาถาด้วยมือจำลอง} สร้างจุดสำคัญของมือสังเคราะห์ที่วาดรูปสามเหลี่ยม หยุดนิ่ง แล้วแบมือ
        ป้อนเข้าเกมจริงแทนผลจากโมเดล เพื่อตรวจว่าวงจร วาด $\rightarrow$ รู้จำ $\rightarrow$ ชาร์จ $\rightarrow$ ปล่อย ทำงานครบ
  \item \textbf{การทดสอบความเท่าเทียมหลังปรับปรุง} เปรียบเทียบผลของโค้ดก่อนและหลังการปรับปรุงประสิทธิภาพ
        ทั้งผลการรู้จำลายเส้นและภาพเอฟเฟกต์
  \item \textbf{การทดสอบประสิทธิภาพ} ตามวิธีในบทที่ 5
\end{itemize}

\section{ผลการทดสอบ}
\begin{itemize}[leftmargin=2cm]
  \item วงจรการร่ายคาถาด้วยมือจำลองทำงานครบ ระบบรู้จำรูปสามเหลี่ยมเป็น Fireball เข้าสู่สถานะชาร์จ
        ใช้เวลาชาร์จ 1.28 วินาที (ค่าที่กำหนดคือ 1.3 วินาที) แล้วสร้างกระสุนได้ถูกต้อง ทั้งเมื่ออัตราการตรวจจับเป็น 30 และ 45 ครั้งต่อวินาที
  \item ตัวรู้จำลายเส้นรุ่นปรับปรุงให้ผลตรงกับรุ่นเดิมทั้ง 400 ลายเส้นที่ทดสอบ
  \item ประสิทธิภาพและระยะการตรวจจับมือเป็นไปตามตารางที่~\ref{tab:fps} และ~\ref{tab:smallhand}
\end{itemize}

\section{ข้อจำกัดของระบบ}
\begin{itemize}[leftmargin=2cm]
  \item \textbf{ยังไม่มีการประเมินความแม่นยำกับผู้ใช้จริงอย่างเป็นระบบ} การทดสอบท่ามือและลายเส้นกับคนจริงทำโดยการสังเกตระหว่างพัฒนา
        ยังไม่ได้เก็บสถิติอัตราการรู้จำถูกต้องจากผู้ใช้หลายคน ส่วนการทดสอบที่มีตัวเลขในรายงานนี้ใช้ข้อมูลสังเคราะห์และภาพนิ่ง
  \item \textbf{การจำแนกท่ามือขึ้นกับมุมของมือ} กฎที่ใช้อาศัยพิกัดสองมิติ เมื่อมือหันข้างหรือชี้เข้าหากล้อง นิ้วจะบังกันและอาจจำแนกผิด
  \item \textbf{ลายเส้นบางคู่คล้ายกัน} จากการทดสอบด้วยลายเส้นสังเคราะห์ที่หมุนและย่อขยายแบบสุ่ม พบว่าวงกลมบางส่วนถูกรู้จำเป็นสี่เหลี่ยม
        และสามเหลี่ยมบางส่วนถูกรู้จำเป็นหัวใจ แม้จะมีการกรองด้วยลักษณะทางเรขาคณิตแล้ว
  \item \textbf{การซูมตรวจจับมือต้องเห็นใบหน้า} หากผู้เล่นหันหลังหรือใบหน้าถูกบัง ระบบจะกลับไปใช้กรอบกว้างซึ่งตรวจจับมือเล็กได้น้อยลง
  \item \textbf{ขึ้นกับแสงและฉากหลัง} โมเดลตรวจจับทำงานได้แย่ลงในที่แสงน้อยหรือเมื่อมีมือของบุคคลอื่นอยู่ในภาพ
  \item \textbf{รองรับเฉพาะ Windows} การเปิดกล้องใช้ DirectShow ซึ่งมีเฉพาะบนระบบปฏิบัติการ Windows
\end{itemize}

% ==============================================================================
\chapter{สรุปผลและข้อเสนอแนะ}
% ==============================================================================

\section{สรุปผล}
โครงงานนี้พัฒนาเกม Magic Spell Duel ได้ตามวัตถุประสงค์ ผู้เล่นสองคนสามารถแข่งขันกันได้ด้วยกล้องเว็บแคมตัวเดียว
โดยร่ายคาถาผ่านการวาดสัญลักษณ์กลางอากาศและการทำท่ามือ ระบบประกอบด้วยการจำแนกท่ามือ 8 ท่า การรู้จำสัญลักษณ์ 6 ชนิด
คาถา 6 ชนิดที่มีสมดุลด้านต้นทุนและผลลัพธ์ กติกาการแข่งขันแบบชนะ 2 ใน 3 ยก และระบบภาพและเสียงครบถ้วน

ในด้านประสิทธิภาพ การแยกงานเป็นหลายเธรดและการปรับอัลกอริทึมทำให้อัตราการแสดงผลเพิ่มขึ้นประมาณ 3--5 เท่า
และขจัดการหยุดชะงักเมื่อรู้จำลายเส้น บทเรียนสำคัญที่ได้คือ ควรวัดก่อนแก้เสมอ
เพราะคอขวดที่แท้จริงหลายจุด (รูปแบบการส่งภาพของกล้อง และตัวรู้จำลายเส้น) ไม่ใช่จุดที่คาดไว้ตั้งแต่แรก
และการเปรียบเทียบผลลัพธ์ก่อนและหลังแก้ไขช่วยให้มั่นใจว่าการปรับปรุงความเร็วไม่ได้เปลี่ยนพฤติกรรมของโปรแกรม

\section{ข้อเสนอแนะและแนวทางพัฒนาต่อ}
\begin{itemize}[leftmargin=2cm]
  \item เก็บข้อมูลการเล่นจากผู้ใช้หลายคนเพื่อวัดอัตราการรู้จำท่ามือและลายเส้น แล้วใช้ปรับค่าเกณฑ์และแม่แบบ
  \item ทดลองใช้ตัวจำแนกท่ามือที่ฝึกจากข้อมูล เพื่อให้ทนต่อมุมของมือได้ดีกว่ากฎเชิงเรขาคณิต
  \item ใช้หน่วยประมวลผลกราฟิกสำหรับโมเดล เพื่อเพิ่มอัตราการตรวจจับมือให้ใกล้เคียงอัตราเฟรมของกล้อง
  \item เพิ่มโหมดผู้เล่นคนเดียวที่ต่อสู้กับคอมพิวเตอร์ และโหมดฝึกวาดสัญลักษณ์
  \item รองรับระบบปฏิบัติการอื่นโดยเลือกวิธีเปิดกล้องตามแพลตฟอร์ม
\end{itemize}

% ==============================================================================
%  บรรณานุกรม
% ==============================================================================
\begin{thebibliography}{9}
\addcontentsline{toc}{chapter}{บรรณานุกรม}

\bibitem{zhang2020hands}
F. Zhang, V. Bazarevsky, A. Vakunov, A. Tkachenka, G. Sung, C.-L. Chang, and M. Grundmann,
``MediaPipe Hands: On-device Real-time Hand Tracking,''
\textit{arXiv preprint arXiv:2006.10214}, 2020.

\bibitem{bazarevsky2019blazeface}
V. Bazarevsky, Y. Kartynnik, A. Vakunov, K. Raveendran, and M. Grundmann,
``BlazeFace: Sub-millisecond Neural Face Detection on Mobile GPUs,''
\textit{arXiv preprint arXiv:1907.05047}, 2019.

\bibitem{wobbrock2007dollar}
J. O. Wobbrock, A. D. Wilson, and Y. Li,
``Gestures without Libraries, Toolkits or Training: A \$1 Recognizer for User Interface Prototypes,''
in \textit{Proceedings of the 20th Annual ACM Symposium on User Interface Software and Technology (UIST '07)}, 2007, pp. 159--168.

\bibitem{lugaresi2019mediapipe}
C. Lugaresi et al.,
``MediaPipe: A Framework for Building Perception Pipelines,''
\textit{arXiv preprint arXiv:1906.08172}, 2019.

\bibitem{bradski2000opencv}
G. Bradski, ``The OpenCV Library,'' \textit{Dr. Dobb's Journal of Software Tools}, 2000.

\bibitem{pygame}
Pygame Community, ``Pygame Documentation,'' \url{https://www.pygame.org/docs/} (สืบค้นเมื่อ [วัน เดือน ปี]).

\end{thebibliography}

% ==============================================================================
%  ภาคผนวก
% ==============================================================================
\appendix

\chapter{คู่มือการติดตั้งและการใช้งาน}

\section{การติดตั้ง}
\begin{lstlisting}[language=bash, numbers=none]
python -m venv venv
venv\Scripts\activate
pip install -r requirements.txt
python game_engine.py
\end{lstlisting}

\section{ปุ่มลัดบนคีย์บอร์ด}
\begin{table}[H]
  \centering
  \caption{ปุ่มลัดบนคีย์บอร์ด}
  \begin{tabular}{l l}
    \toprule
    \textbf{ปุ่ม} & \textbf{การทำงาน} \\
    \midrule
    Space หรือ Enter & เริ่มเกมจากหน้าไตเติล \\
    R & เริ่มการแข่งขันใหม่ \\
    F & สลับโหมดเต็มจอ \\
    A & ปลดล็อกคาถาพิเศษให้ทั้งสองฝั่ง (สำหรับทดสอบ) \\
    Q หรือ Esc & ออกจากเกม \\
    \bottomrule
  \end{tabular}
\end{table}

\section{ค่าตั้งที่ปรับได้บ่อยใน config.py}
\begin{table}[H]
  \centering
  \caption{ค่าตั้งที่สำคัญ}
  \small
  \begin{tabularx}{\textwidth}{l l Y}
    \toprule
    \textbf{ชื่อ} & \textbf{ค่าเริ่มต้น} & \textbf{ความหมาย} \\
    \midrule
    \code{CAMERA_INDEX} & 0 & หมายเลขกล้องที่ใช้ \\
    \code{CAMERA_FPS} & 60 & อัตราเฟรมที่ขอจากกล้อง \\
    \code{MIN_RENDER_FPS} & 75 & อัตราการแสดงผลเป้าหมาย \\
    \code{HAND_MODEL_COMPLEXITY} & 1 & 1 = แม่นยำ, 0 = เร็วกว่าสำหรับเครื่องสมรรถนะต่ำ \\
    \code{HAND_ZOOM_FACE_SCALE} & 10.0 & ขนาดกรอบตรวจจับมือเทียบกับความสูงใบหน้า \\
    \code{MAX_ROUND_WINS} & 2 & จำนวนยกที่ต้องชนะเพื่อชนะการแข่งขัน \\
    \code{STILL_TIMEOUT} & 0.48 & เวลาหยุดนิ้วเพื่อยืนยันการวาด (วินาที, ขั้นต่ำที่ใช้จริงคือ 0.55) \\
    \code{SPELL_CALLOUT_DURATION} & 4.0 & เวลาที่รอท่าปล่อยหลังวาดสำเร็จ (วินาที) \\
    \bottomrule
  \end{tabularx}
\end{table}

\chapter{แหล่งที่มาของทรัพยากรภาพและเสียง}
[ระบุแหล่งที่มาและสัญญาอนุญาตของไฟล์ภาพเอฟเฟกต์ เพลงประกอบ เสียงเอฟเฟกต์ และเสียงผู้ประกาศที่ใช้ในโฟลเดอร์ \code{assets}
รวมถึงส่วนที่ผู้จัดทำสร้างขึ้นเอง]

\end{document}
